\documentclass{paper}

\usepackage[T1]{fontenc}
\usepackage[brazilian]{babel}

\renewcommand{\refname}{Referências}

\title{Aplicação de Conceitos de Sistemas Distribuidos em Aplicações QAOA}
\author{Alexandre Silva \\ \texttt{alexandre.silva1@unesp.br}}

\bibliographystyle{ieeetr}

\begin{document}


\maketitle

\noindent
\textbf{Tema: Como usar métodos e tecnologias clássicas de computação distribuída para resolver QAOA.} 


\section{Introdução}

Computação quântica é uma área de grande interesse nos dias de hoje, visto a teórica capacidade de resolver problemas computacionalmente complexos em tempo hábil. Contudo, computadores quânticos atuais são propensos à erros e não possuem a quantidade necessária de qubits para a implementação de algoritmos considerados mais complexos.

Para solucionar este problema, pesquisadores vêm adotando métodos híbridos, como QAOA. Nestes, um circuito (algoritmo quântico) pequeno é executado em um computador quântico (ou ainda simulador em um computador clássico) e, com base nos resultados, um computador clássico se encarrega de otimizar os parâmetros usados no circuito.

No entanto, mesmo com o emprego de métodos híbridos, a carga necessária tanto de otimização como de execução acabam acarretando no aumento do custo computacional para ambos os lados.

Devido a isto, métodos de distribuição de circuitos entre um conjunto de clusters e QPUs (Quantum Processing Units) tem se destacado como uma solução alternativa para a atual geração de computadores quânticos.

\section{Proposta}

Com isso, neste trabalho, é proposto a revisção dos principais métodos divulgados na literatura para a simulação e execução de QAOA (Quantum Approximate Optimization Algorithm) de forma distribuída (DQAOA) tanto entre clusters como o conjunto cluster + QPUs.


Dentre os artigos usados para a pesquisa preliminar deste trabalho. Encontrou-se a aplicação dos seguintes conhecimentos dentre o leque de possibilidades em computação distribuída:

\begin{enumerate}
	\item Em \cite{GPU} e \cite{QuantumCentric}, foram discutidas algumas maneiras de como repartir um problema baseado em QAOA em sub-QUBOs menores que podem ser processados em nós diferentes dentro de um cluster de forma paralela.
	
	\item Em artigos como \cite{GPU,QuantumCentric,DistributedSimulator,EndToEnd}, são discutidas maneiras de como manter uma comunicação entre nós usando \textit{MPI}. Além disso, também discutem como unir os resultados parciais e como garantir que problemas que possuem \textit{overlap} podem ser resolvidos de forma dsitribuída, garantindo que correlações serão mantidas e a solução global possa ser alcançada.
	
	Outro fator interessante abordado, é como realizar a aplicação de operadores de forma dinâmica entre nós a partir da comunicação entre eles via \textit{MPI} ou similares. Tal método abordado, pode tanto aplicar correções com base em resultados clássicos de medidas, ou ainda a aplicação de denominados \textit{TeleGates}, que são operadores de mais de um qubit que são separados e correlacionados à distância.
	
	\item Também é abordado em \cite{GPU,DistributedSimulator,EndToEnd,Datacenters}, como seria possível utilizar tanto clsuters, dotados de CPU e GPU, assim como QPUs. Nestes casos, também são abordados as camadas de software necessárias para gerenciar estes jobs.
	
	\item Em \cite{Datacenters}, foi mostrado um conjunto de softwares criados com o fim de: repartir o problema, submeter jobs, gerenciar jobs, alocar recursos, unir resultados, entre outros; tudo de forma transparente, utilizando serviços da \textit{AWS} para gerar a infraestrutura necessária. Neste projeto, foram utilizados filas \textit{SQS}, \textit{Elastic} e \textit{Amazon Braket}.
	
\end{enumerate}

Quero abordar estes trabalhos e algumas das técnicas propostas com o enfoque em métodos de computação distribuida.

Além dos exemplos de técnicas citadas acima, existem outras que também foram abordadas nos artigos e quero adicionar neste trabalho, sendo elas:

\begin{enumerate}
	\item Load Balancing;
	\item Caching;
	\item Tolerância à falhas;
\end{enumerate}

Sendo assim, este projeto tem como objetivo entender: \textbf{como usar métodos e tecnologias clássicas de computação distribuída para resolver problemas QAOA}.


\section{Motivação}

A decisão do tema deste trabalho se deu com base no tema da minha pesquisa de mestrado. Nela, utilizarei DQAOA para acelerar certos algoritmos em bioinformática. 

Sendo assim, este trabalho surgiu de encontro com os interesses do meu projeto e me proporcionou a possibilidade de aproveitar o momento para já desbravar a literatura atual.


\bibliography{reference}

\end{document}